\conclusion

В настоящей выпускной квалификационной работе достигнута поставленная цель~--- разработан интеллектуальный программный агент, способный по произвольному словесному описанию сделки определять её гражданско-правовой тип, в диалоговом режиме собирать у пользователя недостающие для подготовки документа сведения и формировать проект договора в формате DOCX вместе с юридическими рекомендациями со ссылками на конкретные статьи нормативных правовых актов. Все поставленные задачи решены, работа выполнена в полном объёме.

Сформирован корпус нормативных правовых актов российского договорного права, включающий часть~первую и часть~вторую Гражданского кодекса Российской Федерации и ряд профильных федеральных законов; реализованы предобработка корпуса, его разбиение на чанки уровня отдельной статьи или её обособленного пункта и индексация в векторной базе данных Qdrant в трёх коллекциях, разнесённых по уровню обобщённости норм. Спроектирована формальная модель агента в виде ориентированного графа состояний из двадцати одного узла, объединённых в семь функциональных фаз обработки запроса; модель реализована программно на фреймворке~LangGraph с поддержкой прерывания обработки на уточнения у пользователя и последующего возобновления с восстановлением промежуточного состояния из контрольного пункта. Реализован механизм диалогового уточнения недостающих сведений: при нехватке данных для согласования существенных условий договора агент формулирует пользователю целевой уточняющий вопрос и продолжает обработку после получения ответа без повторного исполнения предшествующих шагов. Реализована автоматическая итеративная проверка сгенерированного проекта договора на полноту существенных условий и непротиворечивость применимым нормам с возможностью повторной генерации при обнаружении замечаний. Разработан модуль генерации файла договора в формате DOCX с надлежащим оформлением документа делового стиля. Разработаны веб-интерфейс пользователя с двунаправленными уведомлениями через WebSocket и панель администратора для управления корпусом нормативных правовых актов, каталогом поддерживаемых типов сделок и тарифными планами. Обеспечено контейнерное развёртывание программной системы средствами Docker Compose и её работоспособность на локальном оборудовании без передачи пользовательских данных во внешние сервисы. Проведена внешняя экспертная оценка качества подготовленных системой артефактов с привлечением практикующих юристов.

Полученные количественные результаты подтверждают работоспособность принятых проектных решений. Эмпирическая оценка поиска на специально подготовленном датасете из 54 эталонных описаний сделок показала, что в итоговой конфигурации (модель эмбеддингов Giga-Embeddings-instruct, реранкер Qwen3-Reranker-0.6B, LLM-фильтр) достигаются метрики верхней части выдачи \texttt{hit@3}=100~\% и \texttt{precision@3}=91.4~\% при точности классификации типа сделки 100~\%; средняя задержка узла поиска составила 13.6~с, что удерживает полный цикл подготовки документа в рамках поставленного техническим заданием требования к среднему времени отклика. Внешняя экспертная апробация подготовленных артефактов с привлечением трёх практикующих юристов показала, что средняя оценка качества проекта договора по пятибалльной шкале Лайкерта составила 4.42, средняя оценка качества юридических рекомендаций~--- 4.73 при отсутствии оценок ниже уровня <<хорошо>> по обоим показателям и для всех одиннадцати рассмотренных типов сделок. Эти результаты подтверждают практическую применимость предложенного подхода для подготовки договоров по бытовым сделкам гражданского оборота для пользователей без юридического образования.

Научная новизна работы состоит в построении интеллектуального агента, объединяющего в одной программной системе три свойства: работу с произвольной формулировкой сделки от пользователя, опору на проиндексированный корпус российского законодательства и автоматический контроль качества подготавливаемого проекта договора. Подобное сочетание свойств на момент выполнения настоящей работы в открытой литературе и в практических программных продуктах в области правовых технологий не зафиксировано. Практическая значимость работы состоит в том, что разработанная система реализована в виде целостного программного продукта, адаптивного по охвату (перечень поддерживаемых типов сделок настраивается администратором без изменения исходного кода) и сохраняющего конфиденциальность пользовательских данных при локальном развёртывании.

Ограничения проведённого исследования сформулированы открыто и определяют направления дальнейшего развития работы. Целесообразно развёртывание интегральной оценки на несколько раздельных показателей качества (корректность ссылок на нормы права, полнота существенных условий, структурная корректность договора, полнота указания рисков и формальных требований). Дальнейшее развитие системы предполагает расширение каталога поддерживаемых типов сделок за счёт менее распространённых поименованных в Гражданском кодексе Российской Федерации договоров и сопровождение его соответствующим расширением проиндексированного корпуса нормативных правовых актов; расширение библиотеки типовых шаблонов договоров для воспроизведения сложившейся в правоприменительной практике структуры и формулировок документа без отступления от его юридической корректности; интеграцию с государственными информационными системами для автоматического предзаполнения реквизитов сторон по согласию пользователя; развитие поиска за счёт учёта актов толкования права (постановлений и обзоров судебной практики высших судебных инстанций) для повышения качества юридических рекомендаций; горизонтальное масштабирование подсистемы индексации корпуса нормативных правовых актов посредством переноса очереди задач индексации во внешний брокер сообщений (например, Redis) и развёртывания нескольких параллельных воркеров приложения~--- это снимет текущее ограничение, при котором все задачи индексации сериализуются общим внутрипроцессным замком, и позволит одновременно обрабатывать несколько источников при значительном объёме корпуса и при работе нескольких администраторов.
